% 自动生成：python paper/build_tables.py（勿手工编辑）
\begin{table}[htbp]
\centering
\caption{收缩工况代表性时刻的干基含水率（单位：$\mathrm{kg\,kg^{-1}}$）}
\label{tab:t6}
\begin{tabular}{lccccc}
\toprule
时间/h & 0 cm & 0.5 cm & 1 cm & 1.5 cm & 药材表面 \\
\midrule
0 & 2.5500 & 2.5500 & 2.5500 & 2.5500 & 2.5500 \\
6 & 1.7164 & 1.5350 & 1.0212 &  & 0.4215 \\
12 & 0.7340 & 0.6516 & 0.4068 &  & 0.1666 \\
18 & 0.4065 & 0.3667 & 0.2387 &  & 0.0888 \\
24 & 0.2837 & 0.2598 & 0.1783 &  & 0.0671 \\
30 & 0.2253 & 0.2085 & 0.1488 &  & 0.0593 \\
36 & 0.1920 & 0.1790 & 0.1313 &  & 0.0558 \\
42 & 0.1706 & 0.1598 & 0.1196 &  & 0.0539 \\
48 & 0.1556 & 0.1463 & 0.1113 &  & 0.0529 \\
烘干结束时间（50.8235 h） & 0.1500 & 0.1413 & 0.1081 &  & 0.0525 \\
\bottomrule
\end{tabular}
\parbox{0.88\linewidth}{\footnotesize\textit{注：}表中每 6 h 取样；0--1.5 cm 为固定物理坐标，空白表示该位置已位于收缩后药材外部；“药材表面”为移动坐标 $r=R(t)$。末行为连续事件定位得到的精确烘干终止时刻。}
\end{table}
